% ============================================================
% Bayesian Econometrics - Student Presentation I (26.03.2026)
% Paper: Kuschnig & Vashold (2021), BVAR package
% Project: Global Volatility Transmission
% Presenters: Tomas Samaj & Luuk Vos
% ============================================================

\documentclass[10pt, aspectratio=169]{beamer}

% ---------- Theme ----------
\usetheme{metropolis}
\usecolortheme{default}

% ---------- Packages ----------
\usepackage{amsmath, amssymb, bm}
\usepackage{booktabs}
\usepackage{xcolor}
\usepackage[T1]{fontenc}
\usepackage[utf8]{inputenc}
\usepackage[expansion=false]{microtype}
\usepackage{hyperref}

% ---------- Colours ----------
\definecolor{wured}{RGB}{150, 30, 40}
\setbeamercolor{frametitle}{bg=mDarkTeal}
\setbeamercolor{progress bar}{fg=wured}

% ---------- Metadata ----------
\title{BVAR: Bayesian Vector Autoregressions\\with Hierarchical Prior Selection in R}
\subtitle{Kuschnig \& Vashold (2021) \textnormal{\small\itshape Journal of Statistical Software, Vol.~100}}
\author{Tom\'{a}\v{s} \v{S}amaj \and Luuk Vos}
\institute{Bayesian Econometrics \textbar{} WU Vienna}
\date{26 March 2026}

% ============================================================
\begin{document}
% ============================================================

\maketitle

%% ---------------------------------------------------------------
\section{The Paper}
%% ---------------------------------------------------------------

% ------ SLIDE 2: Motivation ------
\begin{frame}{Motivation: The Curse of Dimensionality}

\textbf{Vector Autoregression (VAR)} jointly models $M$ time series over $p$ lags:
\[
  \bm{y}_{t} = \bm{a}_0+ A_1\bm{y}_{t-1} + \cdots + A_p\bm{y}_{t-p} + \bm{\varepsilon}_t,
  \qquad \bm{\varepsilon}_t \sim \mathcal{N}(\bm{0}, \Sigma)
\]

\begin{columns}[T]
  \column{0.52\textwidth}
    \textbf{The problem:}
    \begin{itemize}
      \item Parameter count: $M + M^2 p$ \;\textrightarrow\; grows quadratically
      \item $M{=}6$, $p{=}5$: already \textbf{186 parameters}
      \item OLS overfits; poor out-of-sample forecasting
    \end{itemize}

  \column{0.44\textwidth}
    \textbf{Bayesian solution:}
    \begin{itemize}
      \item Encode prior knowledge via $p(\bm{\theta})$
      \item Regularise via \textbf{shrinkage priors}
      \item Posterior $\propto$ Likelihood $\times$ Prior
    \end{itemize}
\end{columns}

\vspace{0.8em}
\metroset{block=fill}
\begin{block}{Key challenge}
  How to choose the prior's hyperparameters in a principled, data-driven way?
\end{block}

\end{frame}

% ------ INTUITION: VAR ------
\begin{frame}{Intuition: Why Does Dimensionality Matter?}

\begin{columns}[T]
  \column{0.48\textwidth}
    \textbf{The core problem:}
    \begin{itemize}
      \item More variables and lags $\Rightarrow$ exponentially more parameters
      \item OLS must estimate all of them from limited data
      \item The model fits noise rather than signal
    \end{itemize}

  \column{0.48\textwidth}
    \textbf{Why Bayesian shrinkage helps:}
    \begin{itemize}
      \item A prior acts as a soft constraint: it shrinks estimates toward a simple, persistent benchmark
      \item Lets the data override the prior only where evidence is strong
      \item Better out-of-sample forecasts with far fewer ``effective'' parameters
    \end{itemize}
\end{columns}

\vspace{0.8em}
\alert{Key takeaway:} Without regularisation, a VAR with many variables is unreliable. Bayesian priors are a principled way to regularise.

\end{frame}

% ------ SLIDE 3: Minnesota Prior ------
\begin{frame}{The Minnesota Prior}

\textbf{Core idea} (Litterman 1980): each variable is a priori a random walk.
\[
  \mathbb{E}\!\left[(A_s)_{ij}\mid\Sigma\right] =
  \begin{cases} 1 & i = j,\; s = 1 \text{ (first lag)} \\ 0 & \text{otherwise} \end{cases}
\]
\[
  \operatorname{Var}\!\left[(A_s)_{ij}\mid\Sigma\right]
  \propto \frac{\lambda^2 \psi_j}{s^{\,2\alpha}}\cdot\frac{\sigma_{ii}}{\sigma_{jj}}
\]

\vspace{0.5cm}
\begin{columns}[T]
  \column{0.48\textwidth}
    \textbf{Key hyperparameters:}
    \begin{itemize}
      \item $\lambda$ - overall tightness\\
            ($\lambda \to 0$: only prior; $\lambda \to \infty$: only data)
      \item $\alpha$ - lag-decay speed
      \item $\psi_j$ - per-variable scaling
    \end{itemize}

  \column{0.48\textwidth}
    \textbf{Refinements:}
    {\small
    \begin{itemize}
      \item \textbf{Sum-of-coefficients} ($\mu$): shrinks coefficients so persistent variables behave like they stay near their current level
      \item \textbf{Single-unit-root} ($\delta$): shrinks the system toward common stochastic trends, helping accommodate cointegration
    \end{itemize}}
\end{columns}

\vspace{0.5em}
\alert{Open problem:} setting $\lambda, \mu, \delta$ manually is ad hoc and dataset-specific.

\end{frame}

% ------ INTUITION: Minnesota Prior ------
\begin{frame}{Intuition: Why a Random Walk Prior?}


\textbf{The economic logic:}
\begin{itemize}
\item Most macro and financial series are persistent - today's value is a good predictor of tomorrow's
\item Random walk encodes this: ``absent other information, things stay roughly the same''
\item Cross-variable effects are shrunk toward zero - a variable rarely drives another strongly across lags
\end{itemize}

\vspace{0.8em}
\alert{Key takeaway:} The random walk prior is not an assumption we believe literally -- it is a useful regularisation device grounded in stylised facts about time series.

\end{frame}

% ------ SLIDE 4: Hierarchical Prior Selection ------
\begin{frame}{Hierarchical Prior Selection (Giannone et al.\ 2015)}

\textbf{Key idea}: treat hyperparameters $\bm{\gamma} = (\lambda, \mu, \delta, \ldots)$ as random variables.

\bigskip
\textbf{Two-level Bayesian model:}
\[
  p(\bm{\theta}, \bm{\gamma} \mid \bm{y})
  \;\propto\;
  \underbrace{p(\bm{y} \mid \bm{\gamma})}_{\text{marginal likelihood}}
  \; p(\bm{\gamma})
\]
\[
  p(\bm{y} \mid \bm{\gamma})
  = \int p(\bm{y} \mid \bm{\theta}, \bm{\gamma})\; p(\bm{\theta} \mid \bm{\gamma})\; d\bm{\theta}
\]

\begin{itemize}
  \item Marginal likelihood \textbf{penalises complexity} automatically (Occam's razor)
  \item With conjugate NIW prior $\Rightarrow$ \textbf{closed-form} $p(\bm{y}\mid\bm{\gamma})$
  \item $\bm{\gamma}$ drawn via \textbf{Metropolis-Hastings}; $\bm{\theta}$ drawn analytically
\end{itemize}

\vspace{0.4em}
\metroset{block=fill}
\begin{block}{Practical benefit}
  Optimal shrinkage is learned from data - no manual tuning required.
\end{block}

\end{frame}

% ------ INTUITION: Hierarchical Prior Selection ------
\begin{frame}{Intuition: Why Learn Hyperparameters from Data?}

\begin{columns}[T]
  \column{0.48\textwidth}
    \textbf{The problem with manual tuning:}
    \begin{itemize}
      \item Different datasets need different amounts of shrinkage
      \item Short samples: tighter prior helps; long samples: data should dominate
      \item Manually choosing $\lambda$ is arbitrary and hard to justify
    \end{itemize}

  \column{0.48\textwidth}
    \textbf{The hierarchical solution:}
    \begin{itemize}
      \item Marginal likelihood rewards priors that make the data probable
      \item It automatically penalises models that are too complex or too rigid
      \item Can be thought of as ``Bayesian cross-validation'' built into the estimation
    \end{itemize}
\end{columns}

\vspace{0.8em}
\alert{Key takeaway:} Instead of guessing the right amount of shrinkage, the model learns it. This makes results reproducible and principled across different datasets and sample sizes.

\end{frame}

% ------ INTUITION: GFEVD, IRF and Key Outputs ------
\begin{frame}{Intuition: IRF, FEVD and GFEVD}

\textbf{Impulse Response Function (IRF):} traces how a one standard deviation shock to variable $j$ propagates to variable $i$ over time.
\begin{itemize}
  \item Answers: \textit{``If US volatility spikes today, by how much does European volatility rise in 5 days?''}
  \item Recursive (Cholesky) identification assumes a causal ordering; sign-restrictions are an alternative
\end{itemize}

\vspace{0.4em}
\textbf{Forecast Error Variance Decomposition (FEVD):} breaks down the $h$-step forecast uncertainty of variable $i$ into contributions from each structural shock.
\begin{itemize}
  \item Answers: \textit{``What fraction of Japan's volatility forecast error is driven by US shocks?''}
\end{itemize}

\vspace{0.4em}
\textbf{Generalised FEVD (GFEVD)} (Pesaran \& Shin 1998): same idea, but using generalised impulses -- \textbf{no ordering required}.
\begin{itemize}
  \item Robust to variable ordering; especially useful when no clear causal hierarchy exists across markets
\end{itemize}

\end{frame}

% ------ SLIDE 5: The BVAR R Package ------
\begin{frame}{The BVAR R Package}

\begin{columns}[T]
  \column{0.52\textwidth}
    \textbf{Core functions:}
    \begin{enumerate}
      \item \texttt{bv\_minnesota()} - specify prior
      \item \texttt{bvar(data, lags = p, \ldots)} - estimate
      \item \texttt{irf()} / \texttt{predict()} - analyse
    \end{enumerate}

    \vspace{0.5em}
    \textbf{Conjugate Normal-Inverse-Wishart (NIW) setup} (multivariate generalization of the Normal-Inverse-Gamma prior):
    \[
      \bm{\beta}\mid\Sigma \sim \mathcal{N}(\bm{b},\, \Sigma \otimes \Omega)
    \]
    \[
      \Sigma \sim \mathcal{IW}(\Psi, d)
    \]
    $\Rightarrow$ closed-form posteriors; MH only for $\bm{\gamma}$.

  \column{0.44\textwidth}
    \textbf{Features:}
    \begin{itemize}
      \item Recursive \& sign-restricted IRFs
      \item Conditional forecasting (scenario analysis)
      \item FRED-MD / FRED-QD datasets integration
      \item \texttt{coda} interface for convergence diagnostics
      \item \texttt{BVARverse} for \texttt{ggplot2} output
    \end{itemize}
\end{columns}

\end{frame}

%% ---------------------------------------------------------------
\section{Our Project}
%% ---------------------------------------------------------------

% ------ SLIDE 7: Data & Research Question ------
\begin{frame}{Our Project: Data \& Research Question}

\textbf{Research question}: How do volatility shocks transmit across global markets, and does this transmission change during the COVID-19 crisis?

\smallskip
VAR and BVAR applications to macroeconomic and financial variables have proven to be useful tools for studying cross-market dynamics and shock transmission.

\smallskip
\textbf{Data}: daily implied volatility indices, 2015-2025 ($M = 6$, $p=5$): \textbf{186 parameters}

\begin{columns}[T]
  \column{0.54\textwidth}
    \begin{center}
    \small
    \begin{tabular}{lll}
      \toprule
      Index & Market & Underlying \\
      \midrule
      VIX        & United States  & S\&P 500 \\
      V2TX       & Euro Area      & EURO STOXX 50 \\
      IVIUK      & United Kingdom & FTSE 100 \\
      VNKY       & Japan          & Nikkei 225 \\
      VHSI       & Hong Kong      & Hang Seng \\
      India VIX  & India          & Nifty 50 \\
      \bottomrule
    \end{tabular}
    \end{center}

  \column{0.43\textwidth}
    \vspace{0.5cm}
    \textbf{Sub-sample comparison:}
    \begin{itemize}
      \item \textbf{Calm market}: 2015-2019\\(low global volatility)
      \item \textbf{Crisis}: 2020-2021\\(COVID-19 shock)
    \end{itemize}
\end{columns}

\end{frame}

% ------ SLIDE 8: Methodology & Planned Outputs ------
\begin{frame}{Our Project: Methodology \& Planned Outputs}

\begin{columns}[T]
  \column{0.50\textwidth}
    \textbf{Model}: Hierarchical BVAR\\[0.2em]
    \begin{itemize}
      \item Minnesota prior with auto-tuned $\lambda$ (Giannone et al.\ 2015)
      \item Implemented via \texttt{BVAR} R package
      \item Estimated separately on calm and crisis sub-samples
    \end{itemize}

  \column{0.46\textwidth}
    \textbf{Planned outputs:}\\[0.2em]
    \begin{enumerate}
      \item \textbf{IRFs} - 1 s.d.\ US VIX shock; response of all markets over 30-day horizon
      \item \textbf{GFEVD} - share of each market's forecast variance explained by US shocks
      \item \textbf{Sub-sample comparison} - do IRFs \& GFEVDs differ significantly between calm and COVID-19 crisis periods?
    \end{enumerate}
\end{columns}

\end{frame}

% ------ SLIDE 9: References ------
\begin{frame}[allowframebreaks]{References}
\scriptsize
\begin{thebibliography}{99}

\bibitem[Kuschnig \& Vashold, 2021]{Kuschnig2021}
Kuschnig, N., \& Vashold, L. (2021).
\newblock BVAR: Bayesian Vector Autoregressions with Hierarchical Prior Selection in R.
\newblock \emph{Journal of Statistical Software}, 100(14), 1--27.

\bibitem[Giannone et al., 2015]{Giannone2015}
Giannone, D., Lenza, M., \& Primiceri, G. E. (2015).
\newblock Prior Selection for Vector Autoregressions.
\newblock \emph{The Review of Economics and Statistics}, 97(2), 436--451.

\bibitem[Litterman, 1980]{Litterman1980}
Litterman, R. B. (1980).
\newblock Techniques of forecasting using vector autoregressions.
\newblock \emph{Working Paper 115}, Federal Reserve Bank of Minneapolis.

\bibitem[Pesaran \& Shin, 1998]{Pesaran1998}
Pesaran, H. H., \& Shin, Y. (1998).
\newblock Generalized impulse response analysis in linear multivariate models.
\newblock \emph{Economics Letters}, 58(1), 17--29.

\end{thebibliography}
\end{frame}
% ============================================================
\end{document}
% ============================================================
